% =====================================================================
\chapter{Candle Projection: Honest Ghost Candles}
\label{ch:candles}
% =====================================================================

\section{The redesign}

The original design draws a single predicted OHLC sequence. That is
scientifically indefensible and, worse, psychologically dangerous: a crisp
predicted candle invites the user to bet as if it were certain.

The redesign keeps the visual drama and adds truth: draw a \term{probability
cone} made of quantile candles, where the body and wick widths encode
uncertainty and the colour intensity encodes directional confidence.

\begin{center}
\begin{tikzpicture}[scale=0.95]
\draw[msgrey,->] (0,0) -- (10.4,0) node[right,font=\scriptsize]{time};
\draw[msgrey,->] (0,0) -- (0,4.4) node[above,font=\scriptsize]{price};
\node[font=\scriptsize,anchor=north] at (2.5,-0.15) {realised};
\node[font=\scriptsize,anchor=north] at (7.4,-0.15) {projected};
\draw[dashed,msgrey] (5.0,0) -- (5.0,4.3) node[above,font=\scriptsize,black]{now};

% realised candles
\foreach \x/\o/\c/\h/\l in {0.6/1.6/2.0/2.2/1.4, 1.4/2.0/1.8/2.15/1.6,
  2.2/1.8/2.4/2.55/1.7, 3.0/2.4/2.3/2.6/2.15, 3.8/2.3/2.7/2.85/2.2,
  4.5/2.7/2.6/2.9/2.5}{
  \draw[msblue,thick] (\x,\l) -- (\x,\h);
  \draw[msblue,fill=msblue] (\x-0.13,\o) rectangle (\x+0.13,\c);
}
% cone bands
\fill[msblue,opacity=0.10] (5.0,2.6) -- (6.0,3.15) -- (7.0,3.5) -- (8.0,3.75) -- (9.0,3.95) -- (10.0,4.15)
  -- (10.0,0.75) -- (9.0,1.0) -- (8.0,1.3) -- (7.0,1.6) -- (6.0,2.0) -- cycle;
\fill[msblue,opacity=0.18] (5.0,2.6) -- (6.0,2.92) -- (7.0,3.1) -- (8.0,3.25) -- (9.0,3.36) -- (10.0,3.46)
  -- (10.0,1.5) -- (9.0,1.66) -- (8.0,1.82) -- (7.0,2.0) -- (6.0,2.24) -- cycle;
\draw[mscyan,very thick,dashed] (5.0,2.6) -- (6.0,2.56) -- (7.0,2.50) -- (8.0,2.44) -- (9.0,2.39) -- (10.0,2.34);
\node[font=\scriptsize,anchor=west] at (10.05,4.15) {$q_{95}$};
\node[font=\scriptsize,anchor=west] at (10.05,3.46) {$q_{75}$};
\node[font=\scriptsize,anchor=west] at (10.05,2.34) {median};
\node[font=\scriptsize,anchor=west] at (10.05,1.50) {$q_{25}$};
\node[font=\scriptsize,anchor=west] at (10.05,0.75) {$q_{05}$};
% implied move markers
\draw[msamber,thick] (9.6,3.9) -- (10.0,3.9);
\draw[msamber,thick] (9.6,1.1) -- (10.0,1.1);
\node[font=\scriptsize,msamber,anchor=east] at (9.55,3.9) {implied move band};
\end{tikzpicture}
\end{center}

\section{How the cone is produced}

\begin{enumerate}
\item \term{Volatility forecast} $\hat\sigma_h$ from M5 (HAR-RV plus boosted
residual), conditioned on the event.
\item \term{Drift forecast} $\hat\mu_h$ from the calibrated direction ensemble
converted to an expected return: $\hat\mu_h = (2p_{\text{up}}-1)\cdot
\mathbb{E}[|r_h|]$, damped by the meta-labeller's confidence.
\item \term{Direct quantile heads} from M2 give $q_\alpha(r_h)$ for
$\alpha \in \{0.05,0.25,0.5,0.75,0.95\}$ at each horizon.
\item \term{Path sampling} for the visual: draw $N=5000$ paths from a
volatility-clustering process (a student-$t$ innovation GARCH-style simulator, or
a conditional generative model trained on real intraday paths) calibrated so the
terminal distribution matches the quantile heads. Then compute the per-minute
empirical quantiles of open, high, low and close.
\item \term{Rendering}: each projected minute is drawn as a semi-transparent
candle whose body spans $[q_{25}, q_{75}]$ of the close and whose wick spans
$[q_{05}, q_{95}]$, tinted by $P(\text{up})$.
\item \term{Overlay the implied move} from the option chain on the same axis.
When the cone is narrower than the implied move band, options are expensive
relative to our forecast --- that is a sell-volatility setup, and it is
immediately visible.
\end{enumerate}

\begin{keybox}[title={The killer visual}]
Two bands on one chart: \emph{what we forecast} versus \emph{what the market has
priced}. The gap between them is the trade. No competing retail product shows
this, and it makes the product's thesis legible in one glance.
\end{keybox}

\section{Scoring the projection honestly}

\begin{longtable}{@{}p{4.0cm}p{4.6cm}p{5.8cm}@{}}
\toprule
\textbf{Metric} & \textbf{Definition} & \textbf{Target}\\
\midrule
Pinball loss & Quantile loss averaged over $\alpha$ & Must beat a random-walk
baseline with GARCH volatility by at least 8\% relative.\\
CRPS & Continuous ranked probability score & Lower than both random walk and the
implied-volatility-derived distribution.\\
PIT histogram & Probability integral transform of realised outcomes & Should be
uniform. Humps at the edges mean intervals are too narrow.\\
Empirical coverage & Fraction of outcomes inside $[q_{05},q_{95}]$ & Within
$90\% \pm 3\%$.\\
Interval width efficiency & Mean width given coverage & Narrower is better only
if coverage holds.\\
Directional deviation & $|\text{realised} - \text{median}|$ & Reported, but
never as the headline metric.\\
\bottomrule
\end{longtable}

\begin{warnbox}[title={Reject the temptation}]
Do not display a single crisp predicted candle even if users ask for it. Instead
display the median path as a thin dashed line inside the cone. The moment the
product implies certainty, its users will size themselves into ruin and blame
the product --- correctly.
\end{warnbox}

% =====================================================================
\chapter{Options Recommendation Engine}
\label{ch:options}
% =====================================================================

\section{From a price forecast to a trade}

A forecast distribution over the underlying is not a trade. The engine performs
five steps.

\begin{center}
\begin{tikzpicture}[node distance=4mm and 4mm,blk/.append style={text width=2.3cm}]
\node[blk] (a) {1. Fit IV surface\\ SVI / SABR};
\node[blk,right=of a] (b) {2. Forecast\\ $\Delta\sigma_{IV}$ (M6)};
\node[blk,right=of b] (c) {3. Enumerate\\ candidate structures};
\node[blk,right=of c] (d) {4. Monte Carlo\\ P\&L distribution};
\node[blk,right=of d] (e) {5. Rank, size,\\ risk gate};
\draw[ar](a)--(b);\draw[ar](b)--(c);\draw[ar](c)--(d);\draw[ar](d)--(e);
\end{tikzpicture}
\end{center}

\section{Step 1 --- The implied volatility surface}

Raw per-strike implied volatilities from a live chain are noisy and occasionally
arbitrageable due to stale quotes. Fit a smooth, arbitrage-free surface:

\begin{itemize}
\item Compute implied volatility from mid prices only (never last traded price
--- a stale last price generates a fake signal), using Black-76 on futures or
Black--Scholes on spot with the appropriate forward.
\item Filter: bid $>0$, spread below a threshold, open interest above a
threshold, and $|\delta| \in [0.05, 0.95]$.
\item Fit a parametric slice per expiry (SVI is the standard choice), enforcing
no calendar and no butterfly arbitrage.
\item Store surface parameters, not just points. Slope, curvature and level
become features; the fitted surface prices any strike, including illiquid ones.
\end{itemize}

\section{Step 2 --- Forecasting the volatility change}

\begin{keybox}[title={The most valuable model in the system}]
$\Delta\sigma_{IV}$ over the next 15--30 minutes is more forecastable than
direction and it dominates short-dated option P\&L. Typical patterns to learn:
\begin{itemize}
\item Implied volatility \emph{expands} ahead of a scheduled event and
\emph{collapses} within minutes of the release --- regardless of direction.
\item After an unscheduled shock, implied volatility overshoots and mean-reverts
with a half-life measured in tens of minutes.
\item On expiry day, implied volatility decays mechanically and violently.
\item Elevated implied volatility rank means buying premium has a strong
headwind.
\end{itemize}
\end{keybox}

\section{Step 3 --- Candidate structures}

\begin{longtable}{@{}p{3.2cm}p{3.6cm}p{3.4cm}p{3.8cm}@{}}
\toprule
\textbf{Structure} & \textbf{Use when} & \textbf{Main risk} & \textbf{Notes}\\
\midrule
Long ATM/OTM call or put & Strong direction, forecast move well above implied,
implied volatility rank low & Vega and theta & Simple, highest gamma, but the
most exposed to volatility crush.\\
Vertical debit spread & Direction good, implied volatility elevated & Capped
upside & \textbf{Default directional structure post-news.} The short leg
neutralises much of the vega.\\
Long straddle / strangle & Big move expected, direction unclear, before a
scheduled event & Volatility crush if the move disappoints & Only when forecast
move exceeds implied move by a wide margin.\\
Short strangle with wings (iron condor) & Forecast move well below implied, calm
regime & Tail risk & Defined risk only. Never naked short options.\\
Calendar spread & Term structure dislocated & Complex & Advanced; phase 3 of the
options engine.\\
Ratio / broken-wing & Skew dislocation & Asymmetric & Advanced.\\
\textbf{No trade} & Forecast within noise of implied & --- & Must be the most
frequently selected option.\\
\bottomrule
\end{longtable}

Strike selection is constrained, not free: only strikes with adequate open
interest, adequate volume in the last five minutes, and a bid-ask spread below a
fraction of the option price are eligible. A theoretically optimal strike that
cannot be exited is a trap.

\section{Step 4 --- Monte Carlo P\&L distribution}

For each candidate structure:

\begin{enumerate}
\item Take the $N$ sampled underlying paths from Chapter~\ref{ch:candles}.
\item For each path and each future minute, evolve the implied volatility surface
using the M6 forecast plus a stochastic component (correlated negatively with
the underlying for equity indices --- the leverage effect).
\item Reprice the structure along the path, including the passage of time.
\item Apply the exit rule (target, stop, time stop) to get a realised P\&L per
path.
\item Subtract full transaction costs on both legs, entry and exit, including
modelled slippage as a function of the spread and size.
\end{enumerate}

The output is a P\&L distribution, from which we extract:

\[
\mathbb{E}[\text{P\&L}],\quad
P(\text{profit}),\quad
\text{CVaR}_{5\%},\quad
\frac{\mathbb{E}[\text{P\&L}]}{\text{std}},\quad
\text{max drawdown along the path}.
\]

\begin{keybox}[title={Ranking objective}]
Rank by risk-adjusted expected value after costs, penalised by illiquidity:
\[
\text{Score} = \frac{\mathbb{E}[\text{P\&L}] - \lambda\,\text{CVaR}_{5\%}}
{\text{capital at risk}} \cdot \underbrace{\Big(1 - \tfrac{\text{spread}}{2\,\text{mid}}\Big)^2}_{\text{liquidity penalty}}
\]
Emit a recommendation only if $\mathbb{E}[\text{P\&L}] > 0$ after costs
\emph{and} the score exceeds a minimum threshold \emph{and} the meta-labeller
approves.
\end{keybox}

\section{Step 5 --- Position sizing}

Never size by intuition or by a fixed lot count. Use fractional Kelly with a hard
cap:

\[
f^{*} = \frac{p\,b - (1-p)}{b}, \qquad
f_{\text{used}} = \min\!\Big(0.25\,f^{*},\; f_{\max}\Big)
\]

where $p$ is the calibrated win probability, $b$ the payoff ratio from the Monte
Carlo, $0.25$ is the fractional-Kelly discount (full Kelly is far too aggressive
when $p$ is estimated), and $f_{\max}$ is a hard per-trade cap (start at 1--2\%
of capital risked).

\begin{longtable}{@{}p{4.8cm}p{9.8cm}@{}}
\toprule
\textbf{Additional size modifier} & \textbf{Effect}\\
\midrule
Posterior uncertainty of the impact prior & Size scales with
$1/\sqrt{1 + \text{var}}$; small sample sizes shrink the position.\\
Regime & Reduce size in high-volatility or low-liquidity regimes.\\
Time of day & Reduce size in the first 10 and last 15 minutes.\\
Consecutive losses & Reduce size after a losing streak (anti-martingale).\\
Correlation with open positions & Aggregate delta and vega across the book, not
per trade.\\
Expiry day & Halve size; gamma risk is extreme.\\
\bottomrule
\end{longtable}

\section{Complete recommendation record}

Every recommendation is an immutable, fully-explained record. This is both the
product's differentiator and the audit trail.

\begin{lstlisting}[language=Python,caption={Recommendation schema}]
{
  "reco_id": "uuid", "trace_id": "uuid",         # links back to the news item
  "ts_emitted": "2026-08-26T09:20:41.318+05:30",
  "trigger": {"event_id": "...", "class": "geopolitical.military_action",
              "severity_cal": 8.2, "novelty": 0.91, "credibility": 0.88,
              "sources": ["src_a","src_b"], "seconds_since_event": 4.1},
  "forecast": {"horizon_min": 15,
               "p_up": 0.31, "p_up_ci": [0.24, 0.39],
               "quantiles": {"q05": -0.0142, "q25": -0.0071, "q50": -0.0033,
                             "q75": 0.0004, "q95": 0.0061},
               "sigma_forecast": 0.0089, "d_iv_forecast": +2.4,
               "implied_move_15m": 0.0052, "edge_ratio": 1.71},
  "structure": {"type": "bear_put_debit_spread",
                "legs": [{"side":"BUY","strike":19500,"opt":"PE","qty":1},
                         {"side":"SELL","strike":19300,"opt":"PE","qty":1}],
                "net_debit": 78.5, "max_profit": 121.5, "max_loss": 78.5,
                "breakeven": 19421.5},
  "greeks_net": {"delta": -0.21, "gamma": 0.0009, "vega": -1.8, "theta": -3.1},
  "mc": {"paths": 5000, "E_pnl": 21.4, "p_profit": 0.57,
         "cvar5": -71.2, "score": 0.184},
  "sizing": {"kelly_frac": 0.061, "used_frac": 0.015, "lots": 2,
             "capital_at_risk": 11775},
  "exits": {"target": 121.5, "stop": -39.0, "time_stop": "09:36:00",
            "vol_stop": "exit if IV drops > 3 pts"},
  "explanation": ["Military escalation, severity 8.2, first mention 4.1s ago",
                  "Historical prior for this class: median -0.31% over 15m (n=23)",
                  "Forecast |move| 1.71x implied move -> premium buying justified",
                  "IV rank 62 -> spread chosen over naked put to cut vega",
                  "Dealer gamma negative -> moves amplify"],
  "gates_passed": ["credibility>=0.8","two_sources","not_already_priced",
                   "liquidity_ok","daily_loss_limit_ok","regime_ok"],
  "model_versions": {"cls":"ev-cls-v7","dir":"ens-v12","vol":"har-gb-v4",
                     "iv":"ivd-v3","meta":"meta-v5"},
  "mode": "PAPER"
}
\end{lstlisting}

\begin{keybox}[title={Why the explanation field is not decoration}]
Six months from now you will be trying to find out why the system lost money on
a class of events. Without a stored, structured reason for every decision, that
investigation is impossible and the system cannot improve. The explanation array
is a first-class engineering requirement, not a UI nicety.
\end{keybox}

% =====================================================================
\chapter{Risk Management and Kill Switches}
\label{ch:risk}
% =====================================================================

\section{The risk gate sits between the brain and the money}

No component may bypass the risk gate. It is deliberately simple, deliberately
conservative, and deliberately hard to change (changes require a written
justification recorded in the repository).

\begin{longtable}{@{}p{4.4cm}p{3.0cm}p{7.0cm}@{}}
\toprule
\textbf{Limit} & \textbf{Starting value} & \textbf{Behaviour on breach}\\
\midrule
Per-trade capital at risk & 1.5\% & Reject or resize down.\\
Daily loss limit & 3\% & \textbf{Stop for the day.} Flatten everything. No
overrides.\\
Weekly loss limit & 6\% & Stop for the week; mandatory review before restart.\\
Max concurrent positions & 2 & Reject new signals.\\
Net portfolio delta & bounded band & Reject trades that increase concentration.\\
Net portfolio vega & bounded band & Prevents accidental all-in volatility bets.\\
Max trades per day & 6 & Prevents overtrading in a noisy news day.\\
Max trades per event & 1 & Prevents pyramiding on one hypothesis.\\
Slippage guard & 0.5\% of premium & Cancel if fill price deviates.\\
Time-of-day block & first 5 min, last 10 min & No new entries.\\
Expiry-day multiplier & 0.5x size & Gamma risk.\\
Model degraded flag & any & Signals become advisory only.\\
Data staleness & feed silent 3 s & Halt everything.\\
\bottomrule
\end{longtable}

\section{Kill switch hierarchy}

\begin{longtable}{@{}p{2.2cm}p{4.0cm}p{8.4cm}@{}}
\toprule
\textbf{Level} & \textbf{Trigger} & \textbf{Action}\\
\midrule
L1 Soft & Drift alert, calibration decay, unusual latency & Reduce size by half;
notify.\\
L2 Pause & Daily loss limit, data quality failure, conflicting events &
Stop new entries; manage existing positions only.\\
L3 Flatten & Feed loss, model server failure, clock drift, retraction of the
triggering news & Close all positions at market immediately; stop.\\
L4 Lockout & Two L3 events in a week, or any unexplained loss & System refuses to
run until a human clears it with a written root-cause note.\\
\bottomrule
\end{longtable}

\begin{warnbox}[title={Physical kill switch}]
Provide a single command and a single dashboard button that flattens everything
and disables the system, requiring no dependency on the very components that may
have failed. Test it monthly. Also test what happens when the broker API returns
an error during the flatten --- that is the scenario that actually costs money.
\end{warnbox}

\section{Scenario stress tests (run before each promotion)}

\begin{enumerate}
\item Circuit breaker: the index hits a limit and trading halts mid-position.
\item Flash crash and instant recovery within 90 seconds.
\item Broker API returns success but the order was not placed.
\item Duplicate order acknowledgement (double fill).
\item Option chain goes stale but the spot feed stays live.
\item A false catastrophic headline followed by a retraction 4 minutes later.
\item Two contradictory high-severity events 20 seconds apart.
\item Expiry-day final 30 minutes with implied volatility collapsing.
\item Internet loss for 45 seconds with an open position.
\item Model server returns predictions from a stale model version.
\end{enumerate}

For each: define the expected behaviour, implement it, and write an automated
replay test. A system that has not been tested against these has not been
tested.
